\section*{Chapitre \ref{ch:g7:geological-time} --- Les âges de la Terre}

\begin{solution}{exo:g7:geological-time:1}
(a) \omterm{def:g7:geological-time:unconformity}{discordance}~; (b) inclusion~; (c) superposition~; (d) identité
paléontologique~; (e) recoupement.
\end{solution}

\begin{solution}{exo:g7:geological-time:2}
(a) Le grès rouge, en bas, est la plus ancienne~; le grès jaune, en haut,
est la plus récente. (b) Le \omterm{def:g7:geological-time:superposition}{principe de superposition}. (c) Les couches
n'ont pas été retournées depuis leur formation.
\end{solution}

\begin{solution}{exo:g7:geological-time:3}
(a) Paléozoïque, Mésozoïque, Cénozoïque. (b) Le Cénozoïque. (c) Le
Mésozoïque. (d) Le Jurassique vient juste après le Trias~; le Crétacé vient
juste avant le Paléogène.
\end{solution}

\begin{solution}{exo:g7:geological-time:4}
(a) Faux~: la \omterm{def:g7:geological-time:relative-dating}{datation relative} ne fait que ranger les \omterm{def:g3:stones-pebbles-sand:rock}{roches} dans
l'ordre~; un âge en années est un \omterm{def:g7:geological-time:relative-dating}{âge absolu}. (b) Vrai. (c) Faux~: les
\omterm{def:g3:stones-pebbles-sand:pebble}{galets} sont plus anciens que le conglomérat qui les contient (\omterm{def:g7:geological-time:inclusion}{principe
d'inclusion}). (d) Vrai (\omterm{def:g7:geological-time:fossil-succession}{principe d'identité paléontologique}).
\end{solution}

\begin{solution}{exo:g7:geological-time:5}
(a) Plus récent~: il recoupe la couche 3, et une \omterm{def:g3:stones-pebbles-sand:rock}{roche} qui en recoupe une
autre est plus récente (recoupement). (b) Plus ancien~: la couche 4 le
recouvre sans être recoupée. (c) La couche 1, la couche 2, la couche 3, le
granite, la couche 4. (d) Les \omterm{def:g3:stones-pebbles-sand:rock}{roches} qui le recouvraient ont été usées par
l'\omterm{def:g5:water-shapes-the-land:erosion}{érosion} jusqu'à ce que le granite atteigne la surface~; la couche 4
s'est ensuite déposée sur cette surface d'\omterm{def:g5:water-shapes-the-land:erosion}{érosion}, une \omterm{def:g7:geological-time:unconformity}{discordance}.
\end{solution}

\begin{solution}{exo:g7:geological-time:6}
(a) $538.8 - 251.9 = 286.9$~millions d'années. (b) $251.9 - 66 = 185.9$
millions d'années. (c) 66~millions d'années. (d) Le Paléozoïque est la
plus longue, le Cénozoïque la plus courte (pour l'instant).
\end{solution}

\begin{solution}{exo:g7:geological-time:7}
(a) $4540 - 538.8 = 4001.2$~millions d'années. (b) $4001.2 \div 4540
\approx 0.881$~: environ \textbf{88~\%}. (c) Avant le Paléozoïque, les
animaux à coquille n'étaient pas encore courants, et les êtres vivants
mous ne deviennent presque jamais des \omterm{def:g7:sediments-and-fossils:fossil}{fossiles}~; les coquilles et les os,
qui se fossilisent bien, ne sont courants qu'à partir du Paléozoïque.
\end{solution}

\begin{solution}{exo:g7:geological-time:8}
(a) $4540 \div 100 = 45.4$~millions d'années par mètre. (b) $538.8 \div 45.4
\approx 11.9$~: environ \qty{11.9}{m} avant la ligne d'arrivée. (c) $66 \div
45.4 \approx 1.45$~: environ \qty{1.45}{m}. (d) $0.3 \div 45.4 \approx
0.0066$~m, soit environ \qty{0.66}{cm}, moins d'un centimètre~: un coureur
ne le remarquerait guère.
\end{solution}

\begin{solution}{exo:g7:geological-time:9}
(a) Le Silurien est venu d'abord. Oui~: la \omterm{def:g3:stones-pebbles-sand:rock}{roche} la plus ancienne est sous
la plus récente. (b) Les couches grises se déposent à plat~; elles sont
redressées~; elles sont usées~; le \omterm{def:g3:stones-pebbles-sand:sand}{sable} rouge se dépose par-dessus. (c)
C'est une surface d'\omterm{def:g5:water-shapes-the-land:erosion}{érosion} entre des \omterm{def:g3:stones-pebbles-sand:rock}{roches} d'âges très différents~: la
longue durée pendant laquelle les couches grises ont été soulevées et
usées n'a laissé aucune couche.
\end{solution}

\begin{solution}{exo:g7:geological-time:10}
(a) Les deux argilites à ammonites ont le même âge, et les deux calcaires
à oursins aussi~: mêmes \omterm{def:g7:sediments-and-fossils:fossil}{fossiles}, même âge (\omterm{def:g7:geological-time:fossil-succession}{principe d'identité
paléontologique}). (b) La plus ancienne est le grès à trilobites de la
falaise A (il se trouve sous l'argilite à ammonites, qui a le même âge que
la couche du bas de la falaise B). La plus récente est l'argile à feuilles
de la falaise B. (c) De bas en haut~: grès à trilobites~; argilite à
ammonites~; calcaire à oursins~; argile à feuilles.
\end{solution}

\begin{solution}{exo:g7:geological-time:11}
(a) Le schiste. Il se trouve sous le conglomérat (superposition), et les
\omterm{def:g3:stones-pebbles-sand:pebble}{galets} du conglomérat sont des morceaux de schiste (inclusion)~; la
surface d'\omterm{def:g5:water-shapes-the-land:erosion}{érosion} recoupe aussi le schiste. (b) Après~: le grès s'est
déposé à plat, il a donc été incliné plus tard. (c) Avant~: l'argile
recouvre la \omterm{def:g7:earthquakes:fault}{faille} sans être cassée. (d) Le schiste se forme~; il est usé
(surface ondulée)~; le conglomérat, le grès et le calcaire se déposent,
dans cet ordre~; les couches sont inclinées, et la \omterm{def:g7:earthquakes:fault}{faille} les casse ---
la coupe seule ne dit pas lequel de ces deux événements est venu
d'abord~; les couches inclinées et la \omterm{def:g7:earthquakes:fault}{faille} sont usées à plat~; l'argile
se dépose par-dessus. (e) Deux~: la surface ondulée entre le schiste et le
conglomérat, et la surface entre les couches inclinées et l'argile à plat.
\end{solution}

\begin{solution}{exo:g7:geological-time:12}
(a) La moitié de 365~jours, c'est 182.5~jours. De janvier à juin, il y a
$31+28+31+30+31+30 = 181$~jours~: l'histoire est donc à moitié écoulée
1.5~jour après le début de juillet, le \textbf{2~juillet à midi}. (b)
$4540 \div 2 = 2270$~millions d'années. (c) $80 \div \num{4540000000} \times \num{31536000} \approx 0.56$~:
environ \textbf{0.6~seconde}.
\end{solution}

\begin{solution}{pb:g7:geological-time:1}
\textbf{1.} \num{4540}~millions d'années.

\textbf{2.} $4540 \div 365 \approx 12.4$~millions d'années.

\textbf{3.} $\num{4540000000} \div (365 \times 24) = \num{4540000000}
\div 8760 \approx \num{518000}$~ans (avec la valeur arrondie de
\num{12.4}~millions, $\num{12400000} \div 24 \approx \num{517000}$).

\textbf{4.} $\num{4540000000} \div \num{525600} \approx 8638$~: environ
\num{8600}~ans.

\textbf{5.} $4540 - 4030 = 510$~millions d'années après la formation de la
Terre~; $510 \div 12.4 \approx 41$~jours. Janvier compte 31~jours~: 41~jours
se sont donc écoulés au début du \textbf{11~février}.

\textbf{6.} $4540 - 3500 = 1040$~millions d'années~; $1040 \div 12.4 \approx
83.9$~jours (83.6 avec la valeur non arrondie $12.44$).
Janvier et février font 59~jours~: on est 24.6 à 24.9~jours après le
début de mars, le \textbf{25~mars}.

\textbf{7.} $66 \div 12.4 \approx 5.3$~jours avant la fin de l'année,
soit $365 - 5.3 = 359.7$~jours après le 1\textsuperscript{er}~janvier. Les
onze premiers mois font 334~jours~: 25.7~jours après le début de décembre,
le \textbf{26~décembre}.

\textbf{8.} $538.8 \div 4540 \approx 0.119$~: il reste \textbf{11.9~\%} de
l'année, soit $0.119 \times 365 \approx 43$~jours. $365 - 43 =
322$~jours après le 1\textsuperscript{er}~janvier~: il commence en
\textbf{novembre} (le 18~novembre).

\textbf{9.} $\num{10000} \div 8600 \approx 1.2$~minute.

\textbf{10.} $5000 \div 8600 \approx 0.58$~minute, et $0.58 \times 60
\approx 35$~: environ \textbf{35~secondes} avant minuit.

\textbf{11.} $\num{300000} \div 8600 \approx 34.9$~: environ \textbf{35~minutes}
avant minuit.

\textbf{12.} 35~minutes avant minuit le 31~décembre~: notre espèce
apparaît vers \textbf{23~h~25}, dans la dernière heure de l'année.
\end{solution}
